\documentclass[10pt,a4paper]{amsart}
\usepackage[margin=19mm]{geometry}
\usepackage{amsmath,amssymb,mathtools}
\usepackage[scr=rsfs,cal=euler]{mathalfa}
\usepackage{xfrac}
\usepackage[unicode,hidelinks,bookmarks=false]{hyperref}
\newcommand{\ie}{i.e.\ }
\newcommand{\cf}{cf.\ }
\newcommand{\resp}{respectively\ }
\newcommand{\Subsection}{\subsection}
\providecommand{\nobreakdash}{\nobreak\hbox{-}}
\setlength{\emergencystretch}{2em}
\multlinegap0pt
\usepackage{bm}
\usepackage{bbm}
\usepackage{dsfont}
\usepackage{xspace}
\usepackage{cancel}
\usepackage{mathtools}
\usepackage{amsthm}
\makeatletter
\@ifundefined{theorem}{\newtheorem{theorem}{Theorem}}{}
\makeatother
\title[Continuous random variables]{Continuous random variables}
\author{Teo Banica}
\date{Source: arXiv:2608.10415v1 (CC BY 4.0)}
\begin{document}
\maketitle

\noindent\emph{Source and licence.} Continuous random variables. Version arXiv:2608.10415v1 (CC BY 4.0), distributed under the Creative Commons Attribution 4.0 International licence, as recorded in the arXiv licence metadata for this exact version and shown on the abstract page https://arxiv.org/abs/2608.10415v1. Full rights notice: licenses/arxiv-2608.10415-CC-BY-4.0.txt.\par\smallskip

\noindent\textit{Editorial scope.} Counted unit: the Theorem whose source label is \texttt{None} in the article (source0.tex lines 10730--10734, source label \texttt{None}), delivered together with the article's own complete original proof of it (source0.tex lines 10736--10879, one \texttt{proof} environment of 144 lines). No mathematics is written, repaired or re-derived: statement and proof are the article's own text line for line. Required context retained uncounted: none. Every definition, display and label of those lines is retained unchanged. Omitted from this package and not redistributed: 1--10729, 10735--10735, 10880--17013. Nothing inside a retained excerpt is omitted or altered: every retained body is delivered exactly as the article's lines are written, so no omission receipt of any kind accompanies this package. Citations: the retained excerpts carry no citation command at all, so no bibliography entry of the article is transcribed and no reference is redistributed. Reference adaptation: none was needed and none was made; every reference inside the delivered unit resolves inside it, so no unresolved reference or citation is produced. Printed slips kept literal: no printed word, symbol, hypothesis or number of the counted statement or of its proof was altered. Layout and class: the article's own class is \texttt{amsbook} with packages including amssymb, xy, hyperref; the wrapper is the lane's pinned \texttt{amsart} editorial wrapper at 10pt on A4, which restates the theorem-like environments the retained text needs and adds the article's own macro definitions that the retained text needs (none), with extra wrapper packages (bm, bbm, dsfont, xspace, cancel, mathtools). The delivered numbering is the unit's own. Assets: no retained span contains a figure environment, a graphics-inclusion command, a PDF-page inclusion, a table or a TikZ picture; no image file is copied into this package, the package declares no asset dependency and no image file is redistributed with it. Licence: version arXiv:2608.10415v1 of this article is distributed under the Creative Commons Attribution 4.0 International licence, as recorded in the arXiv licence metadata for this exact version and shown on the abstract page https://arxiv.org/abs/2608.10415v1. A full rights notice is delivered at licenses/arxiv-2608.10415-CC-BY-4.0.txt. Characters: every retained excerpt is pure ASCII, so the delivered engine, XeLaTeX with its default Latin Modern text font, reports no missing character.
\begin{theorem}
The density of the chi squared law is given by
$$\chi_n^2=\frac{x^{n/2-1}e^{-x/2}}{2^{n/2}\Gamma(n/2)}\,dx$$
with $\Gamma$ being as usual the gamma function.
\end{theorem}

\begin{proof}
This is something quite tricky, the idea being as follows:

\medskip

(1) To start with, and standing as a useful complement to the above, let us recall from chapter 8 that the gamma function is given by the following formula:
$$\Gamma(s)=\int_0^\infty x^{s-1}e^{-x}\,dx$$

By partial integration we have $\Gamma(s+1)=s\Gamma(s)$, and since $\Gamma(1)=1$, trivially, and $\Gamma(1/2)=\sqrt{\pi}$, obtained via $x=y^2$, we have the following formulae, for $N\in\mathbb N$:
$$\Gamma(N)=(N-1)!\quad,\quad\Gamma\left(N+\frac{1}{2}\right)=\frac{(2N)!!}{2^N}\,\sqrt{\pi}$$

Alternatively, we have the following formula, with $c=\sqrt{2},\sqrt{\pi}$ for $n$ even, odd:
$$\Gamma\left(\frac{n}{2}\right)=\frac{(n-1)!!}{2^{(n-1)/2}}\,c$$

And for more on all this, gamma function and its properties, we refer to chapter 8.

\medskip

(2) Getting now to what our statement says, let us first have a look at the case $n=1$. Here our law is $\chi_1^2=law(f^2)$, with $f\sim g_1$, and the predicted density is:
$$\chi_1^2=\frac{x^{-1/2}e^{-x/2}}{\sqrt{2}\cdot\sqrt{\pi}}\,dx$$

And, is this correct or not? To be more precise, we must prove that we have:
$$f\sim\frac{e^{-x^2/2}}{\sqrt{2\pi}}\,dx\ \implies\ f^2\sim\frac{x^{-1/2}e^{-x/2}}{\sqrt{2\pi}}\,dx$$

In order to do so, we can use moments. Indeed, we have the following formula:
$$M_k(f^2)=M_{2k}(f)=(2k)!!$$

On the other hand, the moments of the predicted density are as follows:
\begin{eqnarray*}
I_k
&=&\frac{1}{\sqrt{2\pi}}\int_0^\infty x^{k-1/2}e^{-x/2}dx\\
&=&\frac{1}{\sqrt{2\pi}}\int_0^\infty\left(k-\frac{1}{2}\right)x^{k-3/2}\cdot 2e^{-x/2}dx\\
&=&(2k-1)\cdot\frac{1}{\sqrt{2\pi}}\int_0^\infty x^{k-3/2}e^{-x/2}dx\\
&=&(2k-1)I_{k-1}
\end{eqnarray*}

As for the initial value, for this recurrence, this is given by the following formula, which by the way shows that we have indeed a probability measure, as we should:
\begin{eqnarray*}
I_0
&=&\frac{1}{\sqrt{2\pi}}\int_0^\infty x^{-1/2}e^{-x/2}dx\\
&=&\frac{1}{\sqrt{2\pi}}\int_0^\infty y^{-1}e^{-y^2/2}\,2y\,dy\\
&=&\sqrt{\frac{2}{\pi}}\int_0^\infty e^{-y^2/2}dy\\
&=&\sqrt{\frac{2}{\pi}}\times\frac{\sqrt{2\pi}}{2}\\
&=&1
\end{eqnarray*}

We conclude that $I_k=(2k)!!$, which matches with $M_k(f^2)=(2k)!!$, so job done.

\medskip

(3) At $n=2$ now, the problematics is closely related to the one for the complex normal laws $G_t$, that we studied in chapter 10, and we can recycle the computations there. To be more precise, with $f,g$ being independent variables, both following $g_1$, we have:
\begin{eqnarray*}
M_k(\chi_2^2)
&=&E((f^2+g^2)^k)\\
&=&\sum_r\binom{k}{r}E(f^{2r})E(g^{2k-2r})\\
&=&\sum_r\binom{k}{r}(2r)!!(2k-2r)!!\\
&=&\sum_r\frac{k!}{r!(k-r)!}\cdot\frac{(2r)!}{2^rr!}\cdot\frac{(2k-2r)!}{2^{k-r}(k-r)!}\\
&=&\frac{k!}{2^k}\sum_r\binom{2r}{r}\binom{2k-2r}{k-r}
\end{eqnarray*}

In order to finish now this moment computation, let us recall from chapter 7 that we have the following formula, coming from the generalized binomial formula:
$$\frac{1}{\sqrt{1+t}}=\sum_{q=0}^\infty\binom{2q}{q}\left(\frac{-t}{4}\right)^q$$

By taking the square of this series, we obtain the following formula:
$$\frac{1}{1+t}
=\sum_k\left(\frac{-t}{4}\right)^k\sum_r\binom{2r}{r}\binom{2k-2r}{k-r}$$

Now by looking at the coefficient of $t^k$ on both sides, we conclude that the sum on the right equals $4^k$. Thus, we can finish our moment computation, as follows:
$$M_k(\chi_2^2)=\frac{k!}{2^k}\times 4^k=2^kk!$$

(4) Still talking $n=2$, let us turn now to the second part of the work which is needed, namely computation of the moments for the predicted density, which is:
$$\chi_2^2=\frac{e^{-x/2}}{2}\,dx$$

As usual by partial integration, the moments of this predicted density are:
\begin{eqnarray*}
I_k
&=&\frac{1}{2}\int_0^\infty x^ke^{-x/2}dx\\
&=&\frac{1}{2}\int_0^\infty kx^{k-1}\cdot 2e^{-x/2}dx\\
&=&2k\cdot\frac{1}{2}\int_0^\infty x^{k-1}e^{-x/2}dx\\
&=&2kI_{k-1}
\end{eqnarray*}

As for the initial value, for this recurrence, this is given by the following formula, which by the way shows that we have indeed a probability measure, as we should:
\begin{eqnarray*}
I_0
&=&\frac{1}{2}\int_0^\infty e^{-x/2}dx\\
&=&\frac{1}{2}\left[-2e^{-x/2}\right]_0^\infty\\
&=&1
\end{eqnarray*}

We conclude that $I_k=2^kk!$, which matches with $M_k(\chi_2^2)=2^kk!$, so job done.

\medskip

(5) After all these preliminaries, time now to get into the general case, $n\in\mathbb N$. With $f_1,\ldots,f_n$ being independent variables, all following the normal law $g_1$, we have:
\begin{eqnarray*}
M_k(\chi_n^2)
&=&E((f_1^2+\ldots+f_n^2)^k)\\
&=&\sum_{i_1+\ldots+i_n=k}\binom{k}{i_1,\ldots,i_n}E(f_1^{2i_1})\ldots E(f_n^{2i_n})\\
&=&\sum_{i_1+\ldots+i_n=k}\frac{k!}{i_1!\ldots i_n!}(2i_1)!!\ldots(2i_n)!!\\
&=&\sum_{i_1+\ldots+i_n=k}\frac{k!}{i_1!\ldots i_n!}\cdot\frac{(2i_1)!}{2^{i_1}i_1!}\cdots\frac{(2i_n)!}{2^{i_n}i_n!}\\
&=&\frac{k!}{2^k}\sum_{i_1+\ldots+i_n=k}\frac{(2i_1)!}{i_1!i_1!}\cdots\frac{(2i_n)!}{i_n!i_n!}\\
&=&\frac{k!}{2^k}\sum_{i_1+\ldots+i_n=k}\binom{2i_1}{i_1}\ldots\binom{2i_n}{i_n}
\end{eqnarray*}

In order to finish this computation, we can use, as in (3), the following formula:
$$\frac{1}{\sqrt{1-4z}}=\sum_{i=0}^\infty\binom{2i}{i}z^i$$

Indeed, by taking the $n$-th power of this series, we obtain the following formula:
$$\frac{1}{\sqrt{(1-4z)^n}}
=\sum_kz^k\sum_{i_1+\ldots+i_n=k}\binom{2i_1}{i_1}\ldots\binom{2i_n}{i_n}$$

Now by putting everything together, the conclusion is that we have:
$$\frac{1}{\sqrt{(1-4z)^n}}
=\sum_{k\geq0}\frac{(2z)^k}{k!}\,M_k(\chi_n^2)$$

But the function on the left can be computed by using the generalized binomial formula, with exponent $p=-n/2$, and we obtain in this way:
\begin{eqnarray*}
\frac{1}{\sqrt{(1-4z)^n}}
&=&\sum_{k\geq0}\binom{-n/2}{k}(-4z)^k\\
&=&\sum_{k\geq0}\frac{(-n/2)(-n/2-1)\ldots(-n/2-k+1)}{k!}(-4z)^k\\
&=&\sum_{k\geq0}\frac{n(n+2)\ldots(n+2k-2)}{k!}(2z)^k
\end{eqnarray*}

And with this, good news, done with our moment computation, which yields:
$$M_k(\chi_n^2)=n(n+2)\ldots(n+2k-2)$$

(6) However, not time to celebrate yet, because we still have to do the second part of the work, namely computation of the moments for the predicted density, which is:
$$\chi_n^2=\frac{x^{n/2-1}e^{-x/2}}{2^{n/2}\Gamma(n/2)}\,dx$$

By using the definition of the gamma function, as formulated in (1), and then several times the formula $\Gamma(s+1)=s\Gamma(s)$, the moments of this predicted density are:
\begin{eqnarray*}
I_k
&=&\frac{1}{2^{n/2}\Gamma(n/2)}\int_0^\infty x^{n/2+k-1}e^{-x/2}\,dx\\
&=&\frac{1}{2^{n/2}\Gamma(n/2)}\int_0^\infty (2y)^{n/2+k-1}e^{-y}\,2dy\\
&=&\frac{2^k}{\Gamma(n/2)}\int_0^\infty y^{n/2+k-1}e^{-y}\,dy\\
&=&2^k\cdot\frac{\Gamma(n/2+k)}{\Gamma(n/2)}\\
&=&2^k\cdot\frac{\Gamma(n/2+k)}{\Gamma(n/2+k-1)}\cdot\frac{\Gamma(n/2+k-1)}{\Gamma(n/2+k-2)}\ldots\frac{\Gamma(n/2+1)}{\Gamma(n/2)}\\
&=&2^k\left(\frac{n}{2}+k-1\right)\left(\frac{n}{2}+k-2\right)\ldots\left(\frac{n}{2}\right)\\
&=&n(n+2)\ldots(n+2k-2)
\end{eqnarray*}

We conclude that we have $I_k=M_k(\chi_n^2)$, so job done, and theorem proved.
\end{proof}

\end{document}
